% Pista pro-25i-q3 · Probabilitat
% Procedència: PAU Catalunya 2025 · Convocatòria de juny (incidències) · Sèrie 4 · Exercici 3
\documentclass[11pt,a4paper]{article}
\usepackage{pista}
\pistainfo{Catalunya 2025}{Convocatòria de juny (incidències)}{Sèrie 4}

\begin{document}

\section*{\textcolor{blauPista}{Exercici 3}}

\noindent Un usuari d'Internet ha estimat que el $25\,\%$ dels correus electrònics que rep són correu brossa, mentre que la resta no ho són. Per a facilitar la classificació del correu, s'ha instal·lat un filtre que envia a la carpeta de correu brossa el $95\,\%$ dels missatges que efectivament ho són. Malauradament, aquest filtre deixa a la safata d'entrada només el $90\,\%$ dels missatges bons (i la resta els envia a la carpeta de correu brossa).

\subsection*{\textit{a)} \normalfont Quina és la probabilitat que un missatge sigui enviat pel filtre a la carpeta de correu brossa? \hfill\textnormal{[0,75 punts]}}

\pista{Un diagrama d'arbre t'ajudarà a visualitzar els dos camins que porten un missatge a acabar a la carpeta de brossa: o bé era efectivament brossa i el filtre l'ha classificat bé, o bé no ho era però el filtre s'ha equivocat. Aplica el teorema de la probabilitat total, i vés amb compte perquè et donen el $90\,\%$ dels missatges bons que \emph{no} van a la brossa: el que et cal és el percentatge que sí que hi va.}

\subsection*{\textit{b)} \normalfont Un dia, aquest usuari obre la carpeta de correu brossa. Quin percentatge de missatges que no són correu brossa hi trobarà? \hfill\textnormal{[0,75 punts]}}

\pista{Aquest és un exemple clar on has d'aplicar el Teorema de Bayes, perquè et demanen una probabilitat condicionada en sentit \emph{invers} respecte de la informació que et donen: tu coneixes la probabilitat de ``anar a la carpeta de brossa \emph{sabent que} no és brossa'', però et demanen la probabilitat de ``no ser brossa \emph{sabent que} és a la carpeta de brossa''.}

\vspace{0.6em}
\noindent En un servidor de correu electrònic determinat, la probabilitat de rebre com a mínim dos correus en menys de $t$ minuts ve donada per la funció
\[
F(t) = \int_{0}^{t} A\,e^{-0{,}5x}\,dx,
\]
on $A$ és una constant real.

\subsection*{\textit{c)} \normalfont Sabent que la probabilitat de rebre com a mínim dos correus en menys de dos minuts és $1 - e^{-1}$, trobeu el valor de $A$. \hfill\textnormal{[1 punt]}}

\pista{Primer resol la integral definida (per trobar una funció primitiva, és bastant directe). Obtindràs $F(t)$ com una expressió que depèn de $t$ i d'$A$. Després, escriu l'equació $F(2) = 1 - e^{-1}$ i aïlla $A$.}

\end{document}
